% Journal of Mathematics and Computer Science submission source.
% Uses the unmodified author-supplied ISRP class dated 2024/08/01.
\pdfminorversion=7
\pdfsuppresswarningpagegroup=1
\documentclass[JMCS,submit]{resources/ISRP}

\usepackage{amsmath,amsthm}
\usepackage{booktabs}
\usepackage{array}
\usepackage{tabularx}
\usepackage{xurl}

\graphicspath{{./}{../../figs/}{figures/}}

\newtheorem{theorem}{Theorem}[section]
\newtheorem{lemma}[theorem]{Lemma}
\newtheorem{proposition}[theorem]{Proposition}
\newtheorem{corollary}[theorem]{Corollary}

\theoremstyle{definition}
\newtheorem{definition}[theorem]{Definition}
\newtheorem{algorithm}[theorem]{Algorithm}
\newtheorem{example}[theorem]{Example}

\theoremstyle{remark}
\newtheorem{remark}[theorem]{Remark}
\numberwithin{equation}{section}

\newcommand{\tas}{\textit{Arabidopsis thaliana}}
\newcommand{\ath}{\textit{A.\ thaliana}}
\newcommand{\hs}{\textit{Homo sapiens}}
\newcommand{\TAU}{\tau}
\newcommand{\xRightarrow}[1]{\stackrel{#1}{\Longrightarrow}}
\newcommand{\LTS}{\mathcal{L}}
\newcommand{\powerset}{\mathcal{P}}
\newcommand{\strongsim}{\mathrel{\sim_{s}}}
\newcommand{\weaksim}{\mathrel{\sim_{w}}}
\newcommand{\weakpre}{\mathrel{\preceq_{w}}}
\newcolumntype{Y}{>{\raggedright\arraybackslash}X}

\hypersetup{
  pdftitle={A graded behavioural comparison framework for labelled Petri nets and qualitative network models},
  pdfauthor={Carlos Ramirez Ovalle},
  pdfkeywords={labelled Petri nets, labelled transition systems, weak bisimulation, simulation preorder, behavioural equivalence, formal verification, model checking}
}

\begin{document}

\title{A graded behavioural comparison framework for labelled Petri nets and
qualitative network models}

\author[a1]{Carlos Ramirez Ovalle\corref{c}}

\address[a1]{Department of Natural Sciences and Mathematics, Pontificia
Universidad Javeriana Cali, Calle 18 No. 118-250, Via a Pance, Cali, Valle del
Cauca, Colombia.}
\authors{C. Ramirez Ovalle}
\cortext[c]{Corresponding author. Email:
\href{mailto:carlosovalle@javerianacali.edu.co}{carlosovalle@javerianacali.edu.co}.}
\copyrightyear{2026}

\begin{abstract}
Given two finite executable labelled models and a declared observable interface,
the comparison problem is to identify the strongest supported relation between
their reachable transition systems without conflating local structure, linear
traces and branching behaviour. We formulate a graded reporting framework whose
exclusive output is strong bisimilarity, weak bisimilarity, mutual weak
simulation, one-way weak simulation or non-comparability. Here ``graded'' denotes
this implication-aware reporting priority, not an algebraic grading. For the
implemented relation-deletion algorithm, we establish monotonicity, finite
termination and greatest-fixed-point correctness. We also state and prove the
relation hierarchy, invariance under the controlled silent edge subdivisions used
by the benchmark, and the inability of label-blind structural comparators to
characterize labelled behaviour in general. Six construction-level cases separate
identity, silent refinement, label mismatch, order change, added behaviour and
trace-equivalent branching. Weak bisimulation classified all 6 cases correctly,
whereas truncated trace equality classified 5, labelled transition-system
graphlet agreement 4, a structural profile 3 and direct PN-GDDA 2. An independent
attacker--defender implementation agreed on all 67,600 ordered comparisons in the
declared universe of one- and two-state systems, and Python and mCRL2 agreed on all
42 reported strong/weak decisions. Public GINsim models, HPN-DREAM/CASPOTS model
families and curated Petri nets demonstrate cross-format use and expose dependence
on update semantics and observational interfaces. These applications support a
reproducible model-level audit; they do not establish predictive performance,
cell-line specificity or organism-level equivalence.

\begin{keyword}
labelled Petri nets \sep labelled transition systems \sep weak bisimulation \sep
simulation preorder \sep behavioural equivalence \sep formal verification \sep
model checking \sep qualitative network models
\MSC{68Q85 \sep 68Q60 \sep 68Q55.}
\end{keyword}
\end{abstract}

\maketitle

\section{Introduction}\label{sec:introduction}

Executable models are compared for several different purposes. An investigator
may ask whether two graphs have similar local structure, whether they generate the
same finite event sequences, whether they preserve branching choices, or whether
the behaviour of one is contained in the other. These questions are not
interchangeable. A structural score may ignore labels and initial states; trace
equality records event order but not the point at which a choice is resolved; and
an equivalence relation does not report asymmetric containment.

Let $N_1$ and $N_2$ be executable labelled models and let $\Sigma$ be a declared
observable alphabet. The problem considered here is
\begin{equation}\label{eq:comparison-map}
 (N_1,N_2,\Sigma)
 \longmapsto (\LTS(N_1),\LTS(N_2))
 \longmapsto C_{\Sigma}(N_1,N_2),
\end{equation}
where $\LTS(N_i)$ is the finite reachable labelled transition system and
$C_{\Sigma}$ returns the strongest supported class among strong bisimilarity,
weak bisimilarity, mutual weak simulation, either one-way weak simulation and
non-comparability. The word \emph{graded} refers to this ordered reporting rule:
relations implied by a stronger result are not reported as separate answers. It
does not denote a graded algebraic structure.

Labelled place/transition nets provide an operational model of causality,
conflict and concurrency \cite{Murata1989}. Their reachable markings induce
labelled transition systems, so standard behavioural relations from concurrency
theory and model checking can be applied \cite{Milner1989,BaierKatoen2008,
Sangiorgi2011}. Internal actions are represented by $\TAU$. Strong bisimulation
matches them as ordinary actions; weak bisimulation and weak simulation abstract
finite internal paths. The resulting distinction is useful whenever two model
descriptions expose the same interface at different internal resolutions.

No single structural or linear-time statistic solves the problem in
\eqref{eq:comparison-map}. Graphlet methods make local structural comparison more
sensitive than elementary size summaries. PN-GDDA, for example, compares
distributions over 151 directed bipartite Petri-net graphlets and 592 published
orbit slots \cite{Szawulak2022}. It intentionally answers a structural question.
In contrast, weak-trace equality is label-aware and global in event order but
forgets branching. The framework therefore keeps exact behavioural relations
separate from diagnostic scores and states which information each output uses.

The formal layer developed below makes the observational interface, simulation
orientation and greatest-fixed-point operator explicit. Its properties justify
the priority used by $C_{\Sigma}$ and show why controlled silent subdivision
preserves weak, but generally not strong, behaviour. Minimal counterexamples
separate equal traces from bisimulation and label-blind structure from labelled
semantics. The implementation is then tested against construction-level ground
truth, mCRL2 and a separately coded attacker--defender oracle. The latter is
exhaustive within the declared finite universe of all one- and two-state systems
over $\{a,\TAU\}$.

Qualitative biological networks provide heterogeneous applications of the formal
problem. Petri nets have long been used for biochemical and regulatory systems
\cite{Koch2005,ChaouiyaPetriBio2007,Heiner2008}, while executable cell biology
requires explicit operational assumptions \cite{Regev2002,Fisher2007}. Boolean
models add a further semantic choice: synchronous, unitary-asynchronous and
most-permissive updates can expose different reachable behaviour
\cite{Pauleve2020}. Model checking has also been applied to Boolean-network
control, while formal constraints support logical-model inference
\cite{CifuentesFontanals2022,Benes2023}. We therefore use externally authored
GINsim models to test cross-format execution, HPN-DREAM/CASPOTS families to
confront formal classes with held-out measurements \cite{Hill2016,Razzaq2018},
and compact curated Petri nets to examine interface-dependent interpretation.
These are applications of the comparison machinery, not premises of its
correctness.

The contributions are:
\begin{enumerate}
\item a formal classifier for two finite labelled models under a declared
observational interface;
\item fixed-point correctness and termination results for the implemented
relation-deletion algorithms, together with the hierarchy that supports the
reporting rule;
\item controlled examples and propositions that separate silent refinement,
label-blind structure, finite traces, branching and directional simulation;
\item a synthetic construction-level benchmark evaluated against trace,
structural-profile, LTS-GDA and direct PN-GDDA comparators;
\item independent checking with mCRL2 and exhaustive agreement with a dual
simulation-game oracle on 67,600 ordered small-LTS pairs;
\item cross-format applications to public GINsim and HPN-DREAM/CASPOTS models and
to curated Petri nets, with update-semantics sensitivity and held-out data kept
distinct from formal correctness; and
\item an end-to-end reproducibility package containing source models, hashes,
tests, an executed notebook and generated outputs.
\end{enumerate}

Table~\ref{tab:novelty} separates the information represented by each comparator.
The integrated audit is the only evaluated path that combines local structure,
global event order, branching, abstraction of hidden events, directional
containment and a held-out data check.

\begin{table}[t]
\caption{Positioning of the evaluated comparison layers. ``Local'' means that
order is represented only inside a bounded motif rather than recursively over
successor states}\label{tab:novelty}
\centering
\scriptsize
\setlength{\tabcolsep}{3pt}
\begin{tabularx}{\textwidth}{@{}>{\raggedright\arraybackslash}Xcccccc@{}}
\toprule
Method & Structure & \shortstack{Global\\order} & Branching &
\shortstack{$\tau$\\abstraction} & \shortstack{Directional\\relation} &
\shortstack{Held-out\\data} \\
\midrule
PN-GDDA \cite{Szawulak2022} & yes & no & no & no & no & no \\
Structural profile & yes & no & no & no & no & no \\
Trace equality & no & yes & no & yes & no & no \\
LTS-GDA & yes/local & no & no & no & no & yes \\
Integrated behavioural audit & yes & yes & yes & yes & yes & yes \\
\bottomrule
\end{tabularx}
\end{table}

The operational sequence is summarised in Fig.~\ref{fig:pipeline}. Biological
delimitation and interface declaration precede formal encoding, reachability
construction and the graded decision rule. This separation makes the modelling
assumptions visible before any equivalence calculation is performed.

\begin{figure}[t]
\centering
\includegraphics[width=0.96\textwidth]{Fig1.pdf}
\caption{Six-stage workflow. Stages 1--3 delimit the model-comparison question
and observational interface; stages 4--6 construct executable Petri nets,
compare their reachable behaviour and assign a model-level relation}
\label{fig:pipeline}
\end{figure}

\section{Mathematical framework}\label{sec:framework}

\subsection{Labelled Petri nets and induced transition systems}
\label{sec:petri}

\begin{definition}[Labelled place/transition net]\label{def:petri-net}
A labelled place/transition net is a tuple
\begin{equation}\label{eq:petri-net}
 N=(P,T,\operatorname{pre},\operatorname{post},\ell,m_0),
\end{equation}
where $P$ and $T$ are finite sets of places and transitions,
$\operatorname{pre},\operatorname{post}:T\to\mathbb{N}^{P}$ are the input and
output multisets, $m_0\in\mathbb{N}^{P}$ is the initial marking, and
$\ell:T\to\Sigma\cup\{\TAU\}$ assigns either an observable action in the declared
interface $\Sigma$ or the internal action $\TAU$.
\end{definition}

Under standard interleaving semantics, an enabled transition $t$ changes marking
$m$ to $m-\operatorname{pre}(t)+\operatorname{post}(t)$. The implementation uses
a worklist to enumerate reachable markings explicitly. A state cap and a token
bound make failure to establish a finite state space explicit rather than silently
truncating the relation calculation. All curated Petri nets used in Section
\ref{sec:applications} are safe, but the definitions below apply to every finite
reachable system.

\begin{definition}[Finite labelled transition system]\label{def:lts}
A finite labelled transition system (LTS) is
\begin{equation}\label{eq:lts}
 L=(S,\Sigma_{\TAU},\longrightarrow,s_0),\qquad
 \Sigma_{\TAU}=\Sigma\cup\{\TAU\},
\end{equation}
where $S$ is a finite state set, $s_0\in S$, and
$\longrightarrow\;\subseteq S\times\Sigma_{\TAU}\times S$. We write
$s\xrightarrow{a}s'$ for $(s,a,s')\in\longrightarrow$. The reachable semantics
of the net in Definition \ref{def:petri-net} is denoted by $\LTS(N)$.
\end{definition}

\subsection{Observable and weak transitions}\label{sec:weak-transitions}

\begin{definition}[Silent closure and weak target sets]\label{def:weak-move}
For $s\in S$, let
\begin{equation}\label{eq:tau-closure}
 \operatorname{Cl}_{\TAU}(s)=
 \{u\in S:s\xrightarrow{\TAU *}u\}.
\end{equation}
For $a\in\Sigma$, define
\begin{equation}\label{eq:weak-target}
 W_a(s)=\{v:\exists u,w\quad
 s\xrightarrow{\TAU *}u\xrightarrow{a}w\xrightarrow{\TAU *}v\},
 \qquad W_{\TAU}(s)=\operatorname{Cl}_{\TAU}(s).
\end{equation}
We write $s\xRightarrow{a}v$ when $v\in W_a(s)$.
\end{definition}

The convention $W_{\TAU}=\operatorname{Cl}_{\TAU}$ is important: a direct silent
move on one side may be matched by zero or more silent moves on the other. It is
also the convention implemented by the Python engine and used in the AUT exports
checked by mCRL2.

\subsection{Simulation and bisimulation}\label{sec:relations}

Let $L_i=(S_i,\Sigma_{\TAU},\longrightarrow_i,s_i^0)$ for $i\in\{1,2\}$.

\begin{definition}[Weak simulation]\label{def:weak-simulation}
A relation $R\subseteq S_1\times S_2$ is a weak simulation from $L_1$ to $L_2$
if, whenever $(x,y)\in R$ and $x\xrightarrow{a}_1x'$, there is
$y'\in W_a^2(y)$ such that $(x',y')\in R$. We write
\begin{equation}\label{eq:simulation-orientation}
 L_1\weakpre L_2
\end{equation}
if $(s_1^0,s_2^0)$ belongs to some such relation. Thus the right-hand system
$L_2$ simulates the left-hand system $L_1$; the left-hand behaviour is contained
in the right-hand behaviour under the interface.
\end{definition}

\begin{definition}[Weak and strong bisimulation]\label{def:bisimulation}
A relation $R\subseteq S_1\times S_2$ is a weak bisimulation if it satisfies the
condition in Definition \ref{def:weak-simulation} and, for every $(x,y)\in R$ and
$y\xrightarrow{a}_2y'$, there is $x'\in W_a^1(x)$ with $(x',y')\in R$. The
initial systems are weakly bisimilar, written $L_1\weaksim L_2$, when their
initial pair belongs to such a relation. Strong bisimulation, written
$L_1\strongsim L_2$, replaces every weak target set $W_a^i(x)$ by the direct
target set $\{x':x\xrightarrow{a}_i x'\}$, including for $a=\TAU$.
\end{definition}

These are standard branching-time relations \cite{Milner1989,Sangiorgi2011}.
The explicit orientation in \eqref{eq:simulation-orientation} prevents the common
ambiguity between ``simulates'' and ``is simulated by'' in one-way results.

\subsection{Weak trace language}\label{sec:trace-language}

\begin{definition}[Exact and truncated weak traces]\label{def:weak-traces}
The finite weak observable language of $L$ is
\begin{equation}\label{eq:exact-language}
 L_w(L)=\{a_1\cdots a_r\in\Sigma^*:s_0
 \xRightarrow{a_1}\cdots\xRightarrow{a_r}s
 \text{ for some }s\in S\}.
\end{equation}
Its depth-$k$ restriction is
$L_{\leq k}(L)=\{u\in L_w(L):|u|\leq k\}$. The diagnostic distance is
\begin{equation}\label{eq:distance}
d_k(L_1,L_2)=1-
\frac{|L_{\leq k}(L_1)\cap L_{\leq k}(L_2)|}
{|L_{\leq k}(L_1)\cup L_{\leq k}(L_2)|}.
\end{equation}
\end{definition}

The implementation computes the prefix-closed truncated sets in Definition
\ref{def:weak-traces}. The statement $d_k=0$ is therefore not identified with
exact weak-trace equivalence. Exact weak-trace results in this study come only
from mCRL2 or from a finite construction for which the complete language is
enumerated analytically.

\subsection{Graded behavioural classification}\label{sec:decision}

\begin{definition}[Graded classifier]\label{def:classifier}
For finite $L_1,L_2$ over the declared interface $\Sigma$, define
$C_{\Sigma}(L_1,L_2)$ by the following priority:
\begin{equation}\label{eq:classifier}
C_{\Sigma}(L_1,L_2)=
\begin{cases}
\mathsf{strong}, & L_1\strongsim L_2,\\
\mathsf{weak}, & L_1\not\strongsim L_2\text{ and }L_1\weaksim L_2,\\
\mathsf{mutual}, & L_1\not\weaksim L_2,\ L_1\weakpre L_2,\ L_2\weakpre L_1,\\
\mathsf{left\_in\_right}, & L_1\weakpre L_2,\ L_2\not\weakpre L_1,\\
\mathsf{right\_in\_left}, & L_2\weakpre L_1,\ L_1\not\weakpre L_2,\\
\mathsf{noncomparable}, & \text{otherwise}.
\end{cases}
\end{equation}
The output is a reporting class, not a newly proposed equivalence relation.
\end{definition}

\begin{proposition}[Exclusive reporting]\label{prop:classifier}
The classifier in Definition \ref{def:classifier} assigns exactly one class to
every ordered pair of finite LTSs.
\end{proposition}

\begin{proof}
The first two cases are selected before any simulation-only case. If neither is
selected, the two simulation predicates have exactly four Boolean combinations:
both true, only the first true, only the second true, or both false. These are the
last four cases of \eqref{eq:classifier}, respectively. The cases are exhaustive
and disjoint by their stated guards. \qed
\end{proof}

\subsection{Fixed-point computation and formal properties}
\label{sec:fixed-point}

For $a\in\Sigma_{\TAU}$, let $D_a^i(x)=\{x':x\xrightarrow{a}_i x'\}$ and
$U=S_1\times S_2$. Define operators on $\powerset(U)$ by
\begin{align}
\Phi_{\preceq}(R)=\{(x,y)\in U:
 &\ \forall a,x'\in D_a^1(x)\ \exists y'\in W_a^2(y),
 (x',y')\in R\},\label{eq:sim-operator}\\
\Phi_w(R)=\{(x,y)\in\Phi_{\preceq}(R):
 &\ \forall a,y'\in D_a^2(y)\ \exists x'\in W_a^1(x),
 (x',y')\in R\},\label{eq:weak-bisim-operator}\\
\Phi_s(R)=\{(x,y)\in U:
 &\ \forall a,x'\in D_a^1(x)\ \exists y'\in D_a^2(y),
 (x',y')\in R,\nonumber\\
 &\ \forall a,y'\in D_a^2(y)\ \exists x'\in D_a^1(x),
 (x',y')\in R\}.\label{eq:strong-bisim-operator}
\end{align}

\begin{theorem}[Correctness of relation deletion]\label{thm:fixed-point}
For each $\Phi\in\{\Phi_{\preceq},\Phi_w,\Phi_s\}$, the operator is monotone.
The implemented algorithm starts with $R=U$, repeatedly deletes any currently
examined pair $z\notin\Phi(R)$, and stops after a complete pass with no deletion.
On finite LTSs it terminates and returns the greatest fixed point $\nu\Phi$.
Consequently, testing membership of $(s_1^0,s_2^0)$ computes, respectively,
weak simulation in the orientation of Definition \ref{def:weak-simulation}, weak
bisimilarity, and strong bisimilarity.
\end{theorem}

\begin{proof}
Let $R\subseteq Q\subseteq U$. Every witness pair required by
\eqref{eq:sim-operator}--\eqref{eq:strong-bisim-operator} that belongs to $R$
also belongs to $Q$. Hence $\Phi(R)\subseteq\Phi(Q)$, proving monotonicity for
all three operators.

Each deletion strictly decreases the finite relation, so at most $|U|$ deletions
can occur. A pass with no deletion therefore occurs after finitely many passes.
Let $G=\nu\Phi$. Initially $G\subseteq U$. Suppose $G\subseteq R$ immediately
before a deletion. If $z\in G$, then
$z\in\Phi(G)\subseteq\Phi(R)$ by fixed-point equality and monotonicity, so $z$
cannot be deleted. Thus $G\subseteq R$ is an invariant.

Let $R_*$ be the relation at termination. Every retained pair passed its transfer
test in the final no-deletion pass, hence $R_*\subseteq\Phi(R_*)$. Conversely,
if $z\notin R_*$, it was deleted from some earlier relation $R_z$ because
$z\notin\Phi(R_z)$. Since $R_*\subseteq R_z$, monotonicity gives
$\Phi(R_*)\subseteq\Phi(R_z)$, and therefore $z\notin\Phi(R_*)$. Thus
$\Phi(R_*)\subseteq R_*$ and $R_*=\Phi(R_*)$.

The invariant gives $G\subseteq R_*$. Since $R_*$ is a fixed point and $G$ is
the greatest fixed point, $R_*\subseteq G$; hence $R_*=G$. Finally,
\eqref{eq:sim-operator}, \eqref{eq:weak-bisim-operator} and
\eqref{eq:strong-bisim-operator} are exactly the transfer clauses in Definitions
\ref{def:weak-simulation} and \ref{def:bisimulation}. \qed
\end{proof}

\begin{remark}[Resource bound]\label{rem:complexity}
Writing $n_i=|S_i|$, the candidate relation uses $O(n_1n_2)$ pairs. At most
$n_1n_2$ deletions and $n_1n_2+1$ full passes are possible. A pass examines the
current pairs and their direct moves; matching scans the corresponding direct or
weak target set. Silent closures are cached, while visible weak target sets are
constructed on demand by the current code. For this reason the implementation
claim is deliberately parameterised by the number of examined pairs, outgoing
moves, weak-target sizes and refinement passes rather than replaced by an
unjustified tight time bound.
\end{remark}

\begin{proposition}[Hierarchy of behavioural relations]\label{prop:hierarchy}
For finite LTSs under Definitions \ref{def:weak-move}--\ref{def:weak-traces},
\begin{equation}\label{eq:hierarchy}
L_1\strongsim L_2
\Longrightarrow L_1\weaksim L_2
\Longrightarrow
(L_1\weakpre L_2\ \land\ L_2\weakpre L_1)
\Longrightarrow L_w(L_1)=L_w(L_2).
\end{equation}
The converses do not hold in general.
\end{proposition}

\begin{proof}
A direct matching move is a permitted weak matching move, so every strong
bisimulation satisfies the weak clauses. Every weak bisimulation satisfies each
one-sided simulation clause by Definition \ref{def:bisimulation}.

It remains to show the language consequence. If $L_1\weakpre L_2$, take a
concrete path witnessing any word in $L_w(L_1)$. Apply the simulation clause to
each direct transition of that path. Observable transitions are matched by
$\TAU^*a\TAU^*$ paths and silent transitions by $\TAU^*$ paths, while successor
pairs remain in the simulation relation. Erasing $\TAU$ therefore gives the same
observable word in $L_2$. Induction on the finite concrete path yields
$L_w(L_1)\subseteq L_w(L_2)$. Simulation in both directions gives equality.
The counterexamples in Proposition \ref{prop:silent-refinement} and Example
\ref{ex:branching} refute the first and last converses. Mutual simulation without
bisimulation is obtained, for example, by comparing a system with initial choices
$a.(b+c)$ and $a.b$ against $a.(b+c)$: each simulates the other, but the extra
$a.b$ successor cannot be related bisimilarly to the state offering both $b$ and
$c$. \qed
\end{proof}

\begin{proposition}[Controlled silent edge refinement]
\label{prop:silent-refinement}
Let $L$ contain $p\xrightarrow{a}q$ with $a\in\Sigma$ and $p\neq q$. Construct
$L'$ by copying all original states and transitions except this edge, and replacing
it by
\begin{equation}\label{eq:silent-refinement}
 \bar p\xrightarrow{a}r\xrightarrow{\TAU}\bar q,
\end{equation}
where $r$ is fresh and has no other incident behaviour. Then $L\weaksim L'$.
Repeated refinements of this form also preserve weak bisimilarity. If $p$ and
$\bar p$ are initial, the selected edge is the only $a$-edge from $\bar p$ in
$L'$, and every $a$-successor of $p$ in $L$ has no outgoing $\TAU$ transition,
then $L\not\strongsim L'$.
\end{proposition}

\begin{proof}
Let
$R=\{(x,\bar x):x\in S\}\cup\{(q,r)\}$. For every copied pair other than the
source of the selected edge, direct transitions are copied and hence can be
matched weakly. At $(p,\bar p)$, the original $a$-move to $q$ is matched by
$\bar p\xrightarrow{a}r\xrightarrow{\TAU}\bar q$, ending in $(q,\bar q)$.
Conversely, $\bar p\xrightarrow{a}r$ is matched by
$p\xrightarrow{a}q$, ending in $(q,r)$.

At $(q,r)$, each move of $q$ is matched after the silent step from $r$ to
$\bar q$. The only direct move from $r$, namely $r\xrightarrow{\TAU}\bar q$,
is matched by the zero-length silent path at $q$, ending in $(q,\bar q)$. Thus
$R$ satisfies both weak transfer clauses. Repeating the construction preserves
weak bisimilarity at each step, so the finite sequence follows by transitivity.

Under the stated additional conditions, a strong match for
$\bar p\xrightarrow{a}r$ must relate $r$ to an $a$-successor $x$ of $p$. The
move $r\xrightarrow{\TAU}\bar q$ would then require a direct $\TAU$-move from
$x$, which is excluded. Hence no strong bisimulation contains the initial pair.
\qed
\end{proof}

\begin{proposition}[Limit of label-blind structural comparison]
\label{prop:label-blind}
Let $M$ be any comparator determined only by unlabelled structure and invariant
under relabelling of transitions. Then $M$ cannot characterize weak-trace
equality, weak simulation or weak bisimilarity for labelled systems in general.
\end{proposition}

\begin{proof}
Take two systems with the same two states and one directed edge. Label the edge
$a$ in $L_a$ and $b\neq a$ in $L_b$. Their unlabelled structures are identical,
so relabelling invariance gives $M(L_a,L_b)=M(L_a,L_a)$. However,
$L_w(L_a)=\{\epsilon,a\}$ and $L_w(L_b)=\{\epsilon,b\}$. Neither initial
observable move can be matched by the other system, so there is no weak simulation
in either direction and no weak bisimulation. Therefore the value assigned to a
self-match by $M$ cannot characterize any of these labelled relations. \qed
\end{proof}

Proposition \ref{prop:label-blind} does not identify a defect in PN-GDDA. It
states that a label-blind structural method answers a different question and
cannot substitute for a labelled behavioural relation.

\begin{example}[Equal traces, different branching]\label{ex:branching}
Let $L_{\mathrm{late}}$ have transitions
\begin{equation}\label{eq:late-choice}
 \ell_0\xrightarrow{a}\ell_1,qquad
 \ell_1\xrightarrow{b}\ell_b,qquad
 \ell_1\xrightarrow{c}\ell_c,
\end{equation}
and let $L_{\mathrm{early}}$ have
\begin{equation}\label{eq:early-choice}
 e_0\xrightarrow{a}e_b,quad e_0\xrightarrow{a}e_c,quad
 e_b\xrightarrow{b}e_b',\quad e_c\xrightarrow{c}e_c'.
\end{equation}
Both exact prefix-closed languages are
$\{\epsilon,a,ab,ac\}$. Nevertheless,
$L_{\mathrm{late}}\not\weaksim L_{\mathrm{early}}$. There are no silent moves,
and matching either early $a$-successor with $\ell_1$ leaves one of the two
late actions unmatched. Directionally,
$L_{\mathrm{early}}\weakpre L_{\mathrm{late}}$: the late-choice state can match
the selected $b$ or $c$ continuation of either early branch. The reverse
$L_{\mathrm{late}}\weakpre L_{\mathrm{early}}$ fails because no single early
$a$-successor offers both continuations. Thus exact trace equality does not
determine branching-time equivalence or simulation orientation.
\end{example}

\subsection{Algorithms}\label{sec:algorithms}

\begin{algorithm}[Reachable LTS construction]\label{alg:reachability}
\textbf{Input:} labelled executable model $N$, initial state or marking $s_0$,
state cap $B$. \textbf{Output:} reachable finite LTS or an explicit cap error.
\begin{enumerate}
\item Initialize a worklist and state index with $s_0$.
\item Remove one discovered state $s$ and enumerate every enabled model move.
\item Execute each move, add its labelled edge, and index any new successor.
\item Abort if the token bound or $B$ is exceeded; otherwise continue until the
worklist is empty.
\end{enumerate}
\end{algorithm}

\begin{algorithm}[Weak targets]\label{alg:weak-targets}
\textbf{Input:} LTS $L$, state $s$, action $a$. \textbf{Output:} $W_a(s)$.
\begin{enumerate}
\item Compute and cache $\operatorname{Cl}_{\TAU}(s)$ by graph search.
\item If $a=\TAU$, return this closure.
\item Otherwise inspect every $a$-edge from every state in the closure and union
the cached silent closures of its targets.
\end{enumerate}
\end{algorithm}

\begin{algorithm}[Greatest-fixed-point relation]\label{alg:fixed-point}
\textbf{Input:} finite $L_1,L_2$ and transfer operator $\Phi$. \textbf{Output:}
whether the initial pair belongs to $\nu\Phi$.
\begin{enumerate}
\item Set $R\leftarrow S_1\times S_2$.
\item Scan a snapshot of the current pairs. Delete $(x,y)$ immediately whenever
$(x,y)\notin\Phi(R)$.
\item Repeat Step 2 until an entire pass deletes no pair.
\item Return whether $(s_1^0,s_2^0)\in R$.
\end{enumerate}
\end{algorithm}

\begin{algorithm}[Graded report]\label{alg:classifier}
\textbf{Input:} finite $L_1,L_2$. \textbf{Output:} one class in Definition
\ref{def:classifier} plus separate diagnostics.
\begin{enumerate}
\item Evaluate strong and weak bisimilarity and both ordered simulations.
\item Apply the priority in \eqref{eq:classifier}; stop at the first satisfied
case.
\item Compute $d_k$ and structural scores only as attached diagnostics.
\end{enumerate}
\end{algorithm}

Table \ref{tab:components} distinguishes exact relations from finite-depth or
structural diagnostics.

\begin{table}[t]
\caption{Inputs and outputs of the computational components}
\label{tab:components}
\centering
\small
\begin{tabularx}{\textwidth}{@{}p{0.22\textwidth}YYp{0.17\textwidth}@{}}
\toprule
Component & Input & Output & Status \\
\midrule
Reachability & labelled executable model & finite reachable LTS & exact under stated caps \\
Strong bisimulation & $L_1,L_2$ & yes/no & formal \\
Weak bisimulation & $L_1,L_2$ & yes/no & formal \\
Weak simulation & ordered $L_1,L_2$ & yes/no & formal \\
Trace distance & $L_1,L_2,k$ & $d_k$ & truncated diagnostic \\
PN-GDDA & Petri-net structure & similarity score & structural diagnostic \\
LTS-GDA & reachable labelled graph & similarity score & local diagnostic \\
\bottomrule
\end{tabularx}
\end{table}

\section{Algorithms and computational validation}
\label{sec:computational-validation}

\subsection{Synthetic benchmark with construction-level ground truth}
\label{sec:benchmark-method}

Six deterministic cases were specified independently of the biological models.
A 12-step labelled workflow is used for the first five cases. The identity case
copies it exactly. Silent refinement subdivides selected observable edges with
internal $\TAU$ steps. Label mismatch replaces a recurring observable. Label
order swap preserves topology and label counts while changing two consecutive
events. Added branch retains the reference workflow and adds one observable
outcome. The sixth case is the standard branching-time distinction
$a.(b+c)$ versus $a.b+a.c$: the trace sets agree, but one model resolves its
choice before the other. Table~\ref{tab:synthetic} records the relation expected
from each construction.

\begin{table}[t]
\caption{Synthetic benchmark and predeclared model-level relation. The observed
class was computed without using the expected label}\label{tab:synthetic}
\centering
\small
\begin{tabularx}{\textwidth}{@{}>{\raggedright\arraybackslash}p{0.18\textwidth}>{\raggedright\arraybackslash}p{0.18\textwidth}YY@{}}
\toprule
Case & Controlled change & Expected class & Observed class \\
\midrule
Identity & exact copy & strong equivalence & strong equivalence \\
Silent refinement & insert internal steps & weak equivalence & weak equivalence \\
Label mismatch & replace one observable family & not comparable & not comparable \\
Label order swap & preserve counts, change order & not comparable & not comparable \\
Added branch & preserve workflow, add outcome & reference simulated by candidate & reference simulated by candidate \\
Trace-equivalent branching & move point of choice & candidate simulated by reference & candidate simulated by reference \\
\bottomrule
\end{tabularx}
\end{table}

\subsection{Independent mCRL2 oracle}\label{sec:mcrl2-method}

To test implementation independence, each benchmark LTS was exported to the
Aldebaran AUT interchange format and evaluated with \texttt{ltscompare} from
mCRL2 202607.0 \cite{Bunte2019}. The external oracle computed strong
bisimilarity, weak bisimilarity and exact weak-trace equivalence, treating
\texttt{tau} as internal. These are separate executions of mCRL2 algorithms on
the exported systems; the Python verdict is not supplied as input. Agreement is
reported for strong and weak bisimulation, while exact weak-trace results are
compared with the trace relation fixed by each construction. The current
\texttt{ltscompare} interface does not provide a weak simulation preorder, so
one-way weak-simulation classifications are not claimed to have an mCRL2 oracle.
They are audited separately with a dual attacker--defender game. Instead of
deleting invalid pairs from a candidate relation, the game grows the least set
of positions from which an attacker can force an unmatched move. Every edge
subset was enumerated for all one- and two-state LTSs over
$\{a,\TAU\}$: 260 systems and $260^2=67{,}600$ ordered comparison pairs. This is
an exhaustive oracle for that declared finite universe, not a proof by testing
for arbitrary model size.

\subsection{Structural and linear-time comparators}\label{sec:baselines}

Four transparent comparators isolate the information contributed by the formal
relation. A structural score averages similarities in state count,
edge count, observable-label multiset and out-degree multiset. Each count
similarity is the smaller-to-larger ratio; multiset similarity is weighted
Jaccard. The trace comparator declares a diagnostic match when $d_8=0$ and
therefore records global order only to the selected depth, with no branching-time
discrimination.

PN-GDDA was implemented directly from its published definition
\cite{Szawulak2022}. All connected induced directed bipartite graphlets on two
through five nodes are generated algorithmically, giving 2, 6, 23 and 120
topologies by size (151 total). For each orbit slot $j$, every Petri-net node's
orbit degree is counted; arc multiplicities, when present, multiply the occurrence
count. If $D_G^j(k)$ is the number of nodes with orbit degree $k$, the normalised
distribution and orbit agreement are
\begin{equation}\label{eq:pn-gdda}
N_G^j(k)=\frac{D_G^j(k)/k}{\sum_{r>0}D_G^j(r)/r},\qquad
A^j(G,H)=1-\frac{1}{\sqrt{2}}
\sqrt{\sum_k\left(N_G^j(k)-N_H^j(k)\right)^2}.
\end{equation}
PN-GDDA is the mean $A^j$ from \eqref{eq:pn-gdda} over the 592 published slots.
The six synthetic LTS pairs are presented to this Petri-net-specific comparator
through an exact state-machine encoding: one place per state, one transition per
LTS edge, one input and one output arc per transition, and one initial token. The
five curated pairs are compared in their original Petri-net form. A score of at
least 0.9 is a predeclared diagnostic decision for the binary benchmark; it is not
treated as a formal equivalence relation.

The catalogue was audited independently against Holmes, the software used for
the published method \cite{Radom2017}. Both hash-recorded Holmes 1.1.1 and
2.0.1.2 contain 151 topologies and 592 orbit slots. A type- and
direction-preserving automorphism reconstruction instead gives 576 distinct
orbits; the discrepancy is confined to 12 five-node graphlets. The primary
analysis preserves the Holmes 592-slot assignment so that the comparator is not
silently redefined, while the 576-orbit partition is reported as a sensitivity
analysis. A reproducible Java probe extracts all 151 topologies and all 592 root
assignments from the hash-verified Holmes 1.1.1 JAR and requires entry-level
agreement with the Python catalogue before evaluating the score. On a four-node
$K_{2,2}$ input net and its one-arc deletion, the independent implementation
reproduced the Holmes score (0.987202746587073 versus 0.9872027465870733), equal
to 12 decimal places.

LTS-GDA is retained as a separate reachable-state comparator. On the simple
undirected LTS projection it evaluates four graphlet orbits (edge, both
three-node-path positions and triangle), then combines their mean agreement
equally with weighted-Jaccard agreement on directed labelled edge, two-step walk,
divergence and convergence motifs. Unlike PN-GDDA, it retains observable labels
and can compare Petri nets, SBML-qual, GINML and inferred Boolean models through
one representation. It uses the same predeclared 0.9 diagnostic threshold.

Binary equivalence ground truth is positive only for identity and silent
refinement. Strong equivalence is included within weak equivalence for this
binary evaluation; simulations are not counted as equivalence. To test whether
the fixed 0.9 result is threshold-dependent, PN-GDDA accuracy is also evaluated
at every distinct decision boundary induced by the six observed scores.

\subsection{Scalability experiment}\label{sec:scalability-method}

Linear workflows of 8, 16, 32, 64, 96 and 128 transitions were compared with an
identical copy and a silently refined version. For each pair, strong bisimulation,
weak bisimulation, both simulation directions and $d_8$ were measured separately.
Wall-clock times are medians of three repetitions. The reference run used Python
3.11.8 on an Apple M3 Pro with 18~GiB memory. Runtime values are machine-dependent;
state counts, edge counts, candidate relation sizes and classifications are
deterministic.

\subsection{Software and reproducibility}\label{sec:repro-method}

The formal engine, public-model parsers and benchmark use the Python standard
library. Plotting and tabulation use NumPy, pandas, Matplotlib and NetworkX
\cite{Hagberg2008};
Petri-net artwork uses pydot and Graphviz. Independent checking uses the
hash-pinned mCRL2 202607.0 release. The PN-GDDA audit records SHA-256 digests of
both inspected Holmes JARs and emits its catalogue and numerical cross-check as
machine-readable JSON. Submission figures are existing vector PDF outputs from
the deterministic analysis pipeline; no figure or result was generated or
modified by a generative-AI system. Fixed seeds control all synthetic
modifications and interface permutations. The repository contains the exact
model registry, benchmark generator, dual simulation-game oracle, unit tests,
executed notebook, environment specifications, continuous-integration workflow
and figure-generation script. Runtime CSV files are excluded from bit-for-bit
checks because wall-clock measurements are hardware-dependent; all other
generated results are deterministic.

\section{Applications to qualitative biological network models}
\label{sec:applications}

\subsection{Controls from public logical models}\label{sec:public-method}

Two externally authored models were obtained from the public GINsim repository
\cite{Naldi2018}: the Boolean mammalian cell-cycle model of Traynard et al.
\cite{Traynard2016}, distributed as SBML-qual \cite{Chaouiya2013SBML}, and the
multilevel p53--Mdm2 DNA-damage model associated with Abou-Jaoud\'e et al.
\cite{AbouJaoude2009}, distributed as GINML. The exact downloaded files, source
URLs and SHA-256 digests are retained in the repository. For the cell-cycle
model, \texttt{CycD=1} and all other components were initially zero. For the
p53--Mdm2 model, the stored GINsim state \texttt{damage} was used.

Both rule sets were interpreted with unitary asynchronous semantics: in one
transition, one enabled non-constant component moves one level toward its
logical target. The observable label records the component and direction
(\texttt{component\_up} or \texttt{component\_down}). Breadth-first enumeration
then constructs the reachable LTS. For each public base model, three controls
were predeclared: an exact copy must be strongly and weakly equivalent; inserting
five internal steps must break strong but preserve weak and weak-trace
equivalence; and renaming one reachable update family must break all three
relations. These controls test parsing, state generation, AUT export and formal
comparison on external dynamics. They do not establish that the published
logical rules are complete biological representations.

As a semantic stress test, each public rule set was also executed synchronously:
all unstable components were evaluated from the current state and moved one
level in a single global transition. A sorted compound label recorded the
component changes. Reachable-state overlap between the asynchronous and
synchronous systems was reported, but no cross-semantics equivalence was
assumed.

\subsection{Blinded HPN-DREAM validation with held-out data}
\label{sec:hpn-method}

The data-backed validation uses the public HPN-DREAM phosphoproteomic time series
for BT20, BT549 and MCF7 breast-cancer cell lines \cite{Hill2016} and the
corresponding CASPOTS Boolean-network families \cite{Razzaq2018}. The archived
families contain 72, 191 and 21 verified
networks, respectively. Files were recovered from public commit
\texttt{95e3c74} of the CASPOTS repository; the prior-knowledge network, learning
sets, model families and mTOR-inhibitor test sets are retained with SHA-256
digests. Learning and test files remain separate throughout the workflow.

One representative per cell line was fixed before reading any test response.
For each family, pairwise Jaccard similarity was computed between active Boolean
clauses. The model with maximum mean within-family similarity was selected; exact
ties were resolved lexicographically on the active-clause tuple using rational
arithmetic. This structure-only rule selected zero-based rows 24, 137 and 9 for
BT20, BT549 and MCF7. An anti-leakage unit test raises an exception if medoid
selection attempts to open a test file.

The observational interface was the intersection of seven measured
PI3K/MAPK/mTOR readouts. It comprised 4EBP1-\allowbreak{}pS65,
AKT-\allowbreak{}pT308, BAD-\allowbreak{}pS112 and
MAPK-\allowbreak{}pT202/\allowbreak{}Y204, together with
MEK1-\allowbreak{}pS217/\allowbreak{}S221, mTOR-\allowbreak{}pS2448 and
p70S6K-\allowbreak{}pT389.
The perturbation set was fixed as
the exact intersection of the three test panels: mTOR inhibition alone,
IGF1 plus mTOR inhibition and 4EBP1 plus mTOR inhibition. Test values determined
neither model selection nor interface membership. Each selected Boolean model
was interpreted with unitary asynchronous semantics. Measured time-zero values
initialized relevant readouts, other relevant nodes were zero, stimuli were
clamped active and inhibitors inactive. Interface updates were labelled by
component and direction; all other updates were $\TAU$. The resulting nine
cross-cell pair--condition comparisons were evaluated by both implementations.

All nine comparisons were then repeated under a global synchronous stress
semantics. Every unstable target was evaluated from the same source state and
updated simultaneously. The transition label recorded the sorted set of changed
interface readouts; a hidden-only global update was $\TAU$. This test asks whether
the categorical relation survives a prominent alternative convention. It does
not implement most-permissive semantics, and the asynchronous analysis remains
primary because CASPOTS evaluates asynchronous trajectories \cite{Pauleve2020,
Razzaq2018}.

Experimental distance was the RMSE between the two cell-line profiles over
common post-zero times and interface readouts. Time-zero measurements initialized
the LTS but were excluded from this outcome to prevent data reuse. Formal classes were ordered from weak
bisimulation through mutual and one-way simulation to non-comparability. Their
association with experimental distance was evaluated with Spearman's $\rho$ and
an exact one-sided test that permuted the three pair labels within each condition
($3!^3=216$ permutations). Truncated trace distance was a secondary metric.
One minus LTS-GDA similarity was evaluated by the same permutation scheme as a
structural comparator; it was not used to select models, conditions or readouts.

Held-out compatibility was also recomputed with CASPOTS/Clingo. The network was
fixed while the solver minimized the published weighted threshold-repair
objective. Because streamed optimization models need not be final, the runner
consumed the solve handle to completion and added deterministic secondary
objectives for continuous squared error and source-row index. Clingo certified
the final optimum, and each of 12 scores was reproduced twice identically. Each
blind medoid was scored on all three test sets; selecting the best network from a
whole family on its own test set was reported separately as a retrospective
oracle, not a blind estimate. Native specificity was tested by exact sign
permutation of the three native-versus-mean-cross differences. Setting every
unmeasured relevant initial node to one was retained as a state-space stress test;
its largest LTS had 27,648 states and was not entered into a quadratic exact
cross-model relation.

\subsection{Illustrative biological model set}\label{sec:case-method}

The application contains five pairs of compact, safe Petri nets: DNA repair and
genome instability (GIM), energy metabolism (DCE), cell-cycle control (SPS), cell
death/autophagy (RCD), and an immune module (AID). Components were selected from
published plant--human comparative studies \cite{Quimbaya2012,
ClavijoBuritica2023} and organized against the cancer-hallmark literature
\cite{Hanahan2011,Hanahan2022}; dedicated reviews informed the qualitative
representation of plant DNA-damage response, cell-cycle control and autophagy
\cite{Yoshiyama2013,DeVeylder2007,Marshall2018}.

Places represent conditions or molecular states, transitions represent curated
events, and the initial marking represents the declared starting context. A
transition is observable only when its label belongs to the cross-system
interface; species-specific intermediate steps are represented as $\TAU$.
Transition-level provenance is stored in the repository and checked for complete
name coverage. These nets are manually curated qualitative abstractions. They do
not encode kinetics, concentrations, stoichiometric uncertainty or tissue
context, and they were not inferred from independent experimental data.

Figure~\ref{fig:gim-petri} exposes the complete GIM encodings rather than only
their computed verdict. Both nets implement damage sensing, checkpoint control,
repair and restoration. Their terminal mechanisms differ: the animal/human net
contains apoptosis, whereas the plant net contains an internal SMR step followed
by differentiation or endoreduplication. The common label \texttt{exit\_cycle} is a
declared coarse observation, not evidence that those terminal mechanisms are
biologically identical.

\begin{figure}[t]
\centering
\begin{minipage}[t]{\textwidth}
  \centering
  \includegraphics[width=0.98\textwidth]{Fig2a.pdf}
  \par\smallskip\textbf{(a) Animal/human GIM model}
\end{minipage}
\par\medskip
\begin{minipage}[t]{\textwidth}
  \centering
  \includegraphics[width=0.98\textwidth]{Fig2b.pdf}
  \par\smallskip\textbf{(b) Arabidopsis GIM model}
\end{minipage}
\caption{Curated DNA-damage Petri nets used in the illustrative comparison:
(a) animal/human and (b) Arabidopsis. Circles denote places and the initial
token; rounded boxes denote transitions.
Observable labels appear in brackets, while grey transitions labelled $\tau$
are internal. The diagrams are generated from the same executable definitions
used by the analysis}
\label{fig:gim-petri}
\end{figure}

The interface is treated as a modelling hypothesis. Diagnostics include targeted
single-label changes, random label permutations, seed sensitivity, a trace-depth
sweep and silent-refinement invariance. The permutation values quantify
sensitivity to scrambled interfaces inside these models; they are not biological
significance tests and do not provide external validation.

\section{Results}\label{sec:results}

\subsection{The synthetic suite recovers every predeclared relation}
\label{sec:benchmark-results}

The formal classification matched the construction-level class in all six cases
(Table~\ref{tab:synthetic}). Identity was strongly bisimilar. Inserting internal
steps broke strong equivalence but preserved weak equivalence. Label mismatch and
label-order change were both rejected. The added branch produced the intended
one-way simulation, while the trace-equivalent branching pair was also separated
by simulation direction rather than incorrectly promoted to equivalence.

For the binary weak-equivalence task, weak bisimulation achieved 6/6 correct
decisions. Trace equality achieved 5/6 because it produced a false positive on
the branching-time case. Direct PN-GDDA achieved 2/6 and classified every pair
as equivalent at the predeclared threshold. It assigned 1.0 to identity, silent
refinement, label mismatch and label-order swap because the first two pairs share
topology and the latter two modify labels that PN-GDDA does not encode. Added
branch and trace-equivalent branching still scored 0.9676 and 0.9657. LTS-GDA
achieved 4/6: silent refinement was a false
negative because graphlets do not abstract $\TAU$, and the added local branch was
a false positive at the predeclared threshold. The structural profile achieved
3/6, with false positives for label mismatch and label-order swap and a false
negative for silent refinement (Table~\ref{tab:baseline};
Fig.~\ref{fig:benchmark}).
No PN-GDDA threshold attained more than 4/6: label mismatch and label-order swap
both tied the two positive controls at 1.0, making those cases structurally
inseparable by any score cutoff.

\begin{table}[t]
\caption{Binary equivalence accuracy on six construction-level cases}
\label{tab:baseline}
\centering
\begin{tabular}{lccc}
\toprule
Method & Accuracy & False positive & False negative \\
\midrule
Weak bisimulation & 1.000 & 0 & 0 \\
Trace equality ($k=8$) & 0.833 & 1 & 0 \\
PN-GDDA-592 ($\geq0.9$) & 0.333 & 4 & 0 \\
LTS-GDA ($\geq0.9$) & 0.667 & 1 & 1 \\
Structural profile ($\geq0.9$) & 0.500 & 2 & 1 \\
\bottomrule
\end{tabular}
\end{table}

\begin{figure}[t]
\centering
\includegraphics[width=\textwidth]{Fig3.pdf}
\caption{Independent benchmark. Left: correctness and predicted binary class for
each case and method. Right: accuracy against the equivalence labels fixed by
construction. The trace baseline misses the branching-time distinction;
PN-GDDA is deliberately label-blind and calls all six encoded nets structurally
similar, whereas labelled LTS-GDA improves on the summary profile but remains
sensitive to silent refinement and local added structure}
\label{fig:benchmark}
\end{figure}

\subsection{Independent software and public-model controls agree}
\label{sec:external-results}

The public cell-cycle model contained 13 Boolean components and generated 826 of
8,192 possible states with 3,413 reachable transitions. The four-component
multilevel p53--Mdm2 model generated 17 of 36 possible states and 26 transitions
(Table~\ref{tab:public-models}). Thus, the external controls exercise the method
on substantially larger and differently encoded state spaces than the compact
curated application.

\begin{table}[t]
\caption{Public logical models used for external workflow validation. Full space
is the product of component level counts; only reachable states enter the formal
comparison}\label{tab:public-models}
\centering
\scriptsize
\setlength{\tabcolsep}{2pt}
\begin{tabularx}{\textwidth}{@{}>{\raggedright\arraybackslash}p{2.25cm}lrrrr>{\raggedright\arraybackslash}X@{}}
\toprule
Model & Format & Comp. & Full space & Reachable & Edges & Initial context \\
\midrule
Mammalian cell cycle & SBML-qual & 13 & 8,192 & 826 & 3,413 & \texttt{CycD=1}; others zero \\
p53--Mdm2 DNA damage & GINML & 4 & 36 & 17 & 26 & stored state \texttt{damage} \\
\bottomrule
\end{tabularx}
\end{table}

All six public-model controls recovered their predeclared relations in both
implementations (Table~\ref{tab:public-controls}). Silent refinement added five
states to each candidate without changing weak behaviour. The observable
perturbations were rejected by strong bisimulation, weak bisimulation and exact
weak-trace equivalence. Across these six controls and the six synthetic pairs,
Python and mCRL2 agreed on 12/12 strong and 12/12 weak decisions. The mCRL2 exact
weak-trace result also matched all 12 predeclared trace controls, including
acceptance of the trace-equivalent branching pair that both bisimulation
implementations reject.

The independent attacker--defender oracle agreed with the production weak-
simulation result on all 67,600 ordered comparisons among the 260 enumerated
small LTSs. Alternative update semantics materially changed public reachability:
the synchronous cell-cycle system contained 13 states, all shared with the 826
asynchronous states (Jaccard 0.0157), while synchronous p53--Mdm2 contained 8
states, all shared with the 17 asynchronous states (Jaccard 0.471). These audits
support implementation correctness while showing that reachable behaviour is
not semantics-neutral.

\begin{table}[t]
\caption{Public-model controls. Entries list strong/weak/weak-trace outcomes;
Python reports the first two relations and mCRL2 reports all three. Y=yes and
N=no}\label{tab:public-controls}
\centering
\scriptsize
\setlength{\tabcolsep}{3pt}
\begin{tabular*}{\textwidth}{@{\extracolsep{\fill}}llccccc}
\toprule
Model & Control & \shortstack{States\\ref./cand.} & \shortstack{Edges\\ref./cand.} & \shortstack{Expected\\S/W/T} & \shortstack{Python\\S/W} & \shortstack{mCRL2\\S/W/T} \\
\midrule
Cell cycle & exact copy & 826/826 & 3,413/3,413 & Y/Y/Y & Y/Y & Y/Y/Y \\
Cell cycle & silent refinement & 826/831 & 3,413/3,418 & N/Y/Y & N/Y & N/Y/Y \\
Cell cycle & observable perturbation & 826/826 & 3,413/3,413 & N/N/N & N/N & N/N/N \\
p53--Mdm2 & exact copy & 17/17 & 26/26 & Y/Y/Y & Y/Y & Y/Y/Y \\
p53--Mdm2 & silent refinement & 17/22 & 26/31 & N/Y/Y & N/Y & N/Y/Y \\
p53--Mdm2 & observable perturbation & 17/17 & 26/26 & N/N/N & N/N & N/N/N \\
\bottomrule
\end{tabular*}
\end{table}

\subsection{Held-out data support compatibility but not prognostic specificity}
\label{sec:hpn-results}

The structure-only representatives contained 23--29 active clauses
(Table~\ref{tab:hpn-medoids}). Across the three shared perturbations, their
reachable LTSs contained 2--44 states. None of the nine cross-cell comparisons
was weakly bisimilar: six showed one-way simulation and three were not comparable.
Python and mCRL2 agreed on every strong and weak decision (18/18), and mCRL2
rejected weak-trace equivalence in all nine pairs.

\begin{table}[t]
\caption{HPN-DREAM model families and blinded structure-only representatives}
\label{tab:hpn-medoids}
\centering
\small
\begin{tabular}{lrrrr}
\toprule
Cell line & Family & Medoid row$^a$ & Clauses & Mean Jaccard \\
\midrule
BT20  & 72  & 24  & 29 & 0.815 \\
BT549 & 191 & 137 & 25 & 0.854 \\
MCF7  & 21  & 9   & 23 & 0.851 \\
\bottomrule
\end{tabular}
\par\smallskip
\raggedright\footnotesize $^a$ Zero-based row; selected from model structure only.
\end{table}

The primary ordinal class had a small positive association with between-cell
experimental distance ($\rho=0.091$), but the exact condition-stratified test was
inconclusive ($p=0.111$). Trace distance was negatively associated with the data
distance ($\rho=-0.596$, $p=0.556$), showing that the truncated linear-time index
is not an empirical similarity surrogate in this small panel. LTS-GDA distance
had the strongest observed association ($\rho=0.669$), but its exact test also
remained inconclusive ($p=0.097$). It is therefore a useful comparator, not a
validated response-distance predictor.

Under synchronous updating, four pairs had one-way simulation and five were not
comparable. Seven of nine asynchronous classes were preserved. The
BT20--BT549 pair under IGF1 plus mTOR inhibition and the BT20--MCF7 pair under
4EBP1 plus mTOR inhibition changed from one-way simulation to
non-comparability. Thus the broad absence of equivalence persisted, but two
directional-containment statements were update-semantics dependent.

CASPOTS reproduced compatibility with the held-out test data. Native medoid RMSE
was 0.33024 for BT20, 0.31335 for BT549 and 0.27726 for MCF7; corresponding
discretization RMSE was 0.32944, 0.30079 and 0.27726. The native medoid was best
or tied for BT20 and MCF7, but not for BT549 (Fig.~\ref{fig:hpn-validation}a).
Across the three targets, the native excess-RMSE advantage over the mean of both
cross-cell medoids was 0.00212 and was not significant under the exact paired
sign test ($p=0.25$). The independently selected formal class likewise did not
significantly order experimental distances (Fig.~\ref{fig:hpn-validation}b).
These results validate held-out compatibility and deterministic reproduction of
the published scoring logic, while directly rejecting a claim of demonstrated
cell-specific predictive discrimination.

\begin{figure}[t]
\centering
\includegraphics[width=\textwidth]{Fig8.pdf}
\caption{Blinded HPN-DREAM validation. (a) CASPOTS excess RMSE above the
discretization reference for every structure-only source medoid and held-out
target cell line; blue boxes mark native source--target combinations. (b)
Between-cell profile RMSE under the primary formal class for three shared
perturbations. Horizontal lines are class medians. The exact test permutes cell
pairs within perturbation; the annotation reports both ordinal-class and
LTS-GDA-distance associations. Neither panel supports robust native-cell
specificity}
\label{fig:hpn-validation}
\end{figure}

\subsection{Runtime grows with the candidate relation}\label{sec:scaling-results}

Across the two scaling series, complete runtime rose from approximately
0.4~ms for 9-state reference models to approximately 0.54~s for the largest
pairs. The 128-step identity pair contained $129\times129=16{,}641$ candidate
state pairs; silent refinement increased the candidate to 145 states and the
relation to 18,705 pairs. Both series retained their expected classifications at
every size (Fig.~\ref{fig:scaling}). Two simulation computations accounted for
the largest runtime share. These measurements establish feasibility for compact
and medium qualitative modules, not for genome-scale state spaces.

\begin{figure}[t]
\centering
\includegraphics[width=0.96\textwidth]{Fig4.pdf}
\caption{Scalability on deterministic workflows. Left: median total runtime
against candidate state-pair count on logarithmic axes. Right: decomposition of
the largest run into strong bisimulation, weak bisimulation, two simulations and
trace distance. Timings are machine-dependent; model sizes and outcomes are
deterministic}
\label{fig:scaling}
\end{figure}

\subsection{Case-study outputs are conditional relations between curated models}
\label{sec:case-results}

The ten curated Petri nets generated reachable systems of 5--14 states. Under the
declared interfaces, the GIM, DCE and SPS model pairs were weakly, but not
strongly, bisimilar and had $d_6=0$. The RCD plant model was simulated by the
animal model but not conversely ($d_6=0.143$). The AID pair satisfied neither
simulation direction ($d_6=0.429$; Table~\ref{tab:case}).

\begin{table}[t]
\caption{Relations between the curated case-study model pairs. These are
model-level, interface-conditional outputs; $P$ denotes the plant model and
$A$ the animal model}\label{tab:case}
\centering
\small
\begin{tabular*}{\textwidth}{@{\extracolsep\fill}llccccc}
\toprule
Module & Encoded process & Strong & Weak & $P\preceq A$ & $A\preceq P$ & $d_6$ \\
\midrule
GIM & DNA repair & no & yes & yes & yes & 0.000 \\
DCE & energy metabolism & no & yes & yes & yes & 0.000 \\
SPS & cell-cycle control & no & yes & yes & yes & 0.000 \\
RCD & cell death/autophagy & no & no & yes & no & 0.143 \\
AID & immune control & no & no & no & no & 0.429 \\
\bottomrule
\end{tabular*}
\end{table}

Direct PN-GDDA reached the same binary decision for all five native pairs despite
their different behavioural classes (Table~\ref{tab:pn-gdda-native}). The RCD
nets received a perfect structural score although only one simulation direction
held; AID scored 0.9907 although neither direction held. Replacing the 592 Holmes
slots with the 576 automorphism-distinct orbits changed no score by as much as
0.001 and changed no threshold decision. Thus, the disagreement is not caused by
the catalog audit: local unlabelled Petri-net structure and reachable labelled
behaviour answer different questions.

\begin{table}[t]
\caption{Direct PN-GDDA on the original curated Petri-net pairs. ``592'' is the
published/Holmes slot catalog; ``576'' is the automorphism-partition sensitivity.
The $\geq0.9$ column is a diagnostic threshold, not a formal relation}
\label{tab:pn-gdda-native}
\centering
\small
\begin{tabular*}{\textwidth}{@{\extracolsep\fill}lcccl}
\toprule
Module & PN-GDDA-592 & Orbit-576 & $\geq0.9$ & Behavioural relation \\
\midrule
GIM & 0.9976 & 0.9975 & yes & weak bisimulation \\
DCE & 0.9969 & 0.9968 & yes & weak bisimulation \\
SPS & 1.0000 & 1.0000 & yes & weak bisimulation \\
RCD & 1.0000 & 1.0000 & yes & one-way simulation \\
AID & 0.9907 & 0.9904 & yes & not comparable \\
\bottomrule
\end{tabular*}
\end{table}

In GIM, apoptosis is the animal terminal fate, whereas differentiation and
endoreduplication are plant terminal fates. All three outcomes are represented
through the common coarse observable \texttt{exit\_cycle}. The weak-equivalence result therefore says that
the two encoded models agree at this selected resolution. A finer interface that
distinguishes those outcomes would change the classification. This example makes
the conditional nature of the method explicit rather than treating the interface
as a neutral measurement device.

The corresponding reachability systems in Fig.~\ref{fig:gim-lts} make the weak
relation inspectable. Observable paths for damage, checkpoint, repair,
restoration and cycle exit can be matched, while additional internal paths are
absorbed by $\TAU$ closure. The figure therefore connects the Petri-net
construction to the state-pair relation used by the algorithm.

\begin{figure}[t]
\centering
\begin{minipage}[t]{0.49\textwidth}
  \centering
  \includegraphics[width=\textwidth]{Fig5a.pdf}
  \par\smallskip\textbf{(a) Animal/human reachability system}
\end{minipage}\hfill
\begin{minipage}[t]{0.49\textwidth}
  \centering
  \includegraphics[width=\textwidth]{Fig5b.pdf}
  \par\smallskip\textbf{(b) Arabidopsis reachability system}
\end{minipage}
\caption{Reachability systems generated from the GIM Petri nets: (a)
animal/human and (b) Arabidopsis. Numbered nodes are reachable markings, the
initial node is distinguished by its fill, solid
labelled arcs are observable moves and dashed arcs are internal $\tau$ moves}
\label{fig:gim-lts}
\end{figure}

The RCD pair supplies a contrasting one-way result (Fig.~\ref{fig:rcd-petri}).
Both encodings include stress, mTOR inhibition, autophagy and survival. The
animal model additionally exposes canonical caspase apoptosis, whereas the plant
model proceeds through an internal metacaspase event. Consequently, the plant
behaviour is simulated by the animal model under the declared interface, but the
reverse direction fails. This is a containment statement about the two encodings,
not a claim that either net exhausts cell-death biology.

\begin{figure}[t]
\centering
\begin{minipage}[t]{\textwidth}
  \centering
  \includegraphics[width=0.94\textwidth]{Fig6a.pdf}
  \par\smallskip\textbf{(a) Animal RCD model with an observable caspase-apoptosis branch}
\end{minipage}
\par\medskip
\begin{minipage}[t]{\textwidth}
  \centering
  \includegraphics[width=0.94\textwidth]{Fig6b.pdf}
  \par\smallskip\textbf{(b) Arabidopsis RCD model with an internal metacaspase branch}
\end{minipage}
\caption{Curated cell-death/autophagy Petri nets illustrating the source of the
one-way simulation result: (a) animal and (b) Arabidopsis. Symbols and transition styling follow
Fig.~\ref{fig:gim-petri}}
\label{fig:rcd-petri}
\end{figure}

Randomly subdividing internal structure preserved every categorical result, as
required by weak semantics. Renaming individual observables changed the expected
model relations, confirming that interface labels affect the computed result.
Label-permutation diagnostics separated the three weakly equivalent curated
pairs from scrambled versions, but these tests only assess the internal
specificity of the selected interface. They cannot establish biological validity
because the original labels and networks were manually curated
(Fig.~\ref{fig:diagnostics}).

\begin{figure}[t]
\centering
\begin{minipage}[t]{0.47\textwidth}
  \centering
  \includegraphics[width=\textwidth]{Fig7a.pdf}
  \par\smallskip\textbf{(a) Silent-refinement invariance}
\end{minipage}\hfill
\begin{minipage}[t]{0.51\textwidth}
  \centering
  \includegraphics[width=\textwidth]{Fig7b.pdf}
  \par\smallskip\textbf{(b) Scrambled-interface reference}
\end{minipage}
\caption{Internal robustness diagnostics for the curated case study: (a)
verdict preservation under 300 silent refinements per module and (b) observed
distance against 2,000 scrambled-interface distances. Bars and numeric labels
report refinement invariance; violin distributions show the label-permutation
reference and points show observed distances. Reported
$p$- and Benjamini--Hochberg $q$-values quantify model-internal interface
specificity and are not experimental biological significance tests}
\label{fig:diagnostics}
\end{figure}

\section{Discussion}\label{sec:discussion}

\subsection{Semantic meaning of the graded comparison}
\label{sec:discussion-meaning}

The classifier turns an ordered pair of executable models into one
interface-conditional statement. Its priority is supported by Proposition
\ref{prop:hierarchy}: a strong result already entails the corresponding weak and
simulation results, while a weak result already entails both simulation
directions. Reporting only the strongest satisfied class prevents the same pair
from receiving several apparently competing labels. The one-way classes preserve
orientation and therefore express containment rather than equivalence.

The fixed-point result is about the algorithm that was executed, not an idealized
variant. Theorem \ref{thm:fixed-point} covers the in-place deletion order, the
implemented treatment of $\TAU$, and termination on finite state spaces. It also
separates mathematical correctness from software validation: the proof applies
to arbitrary finite inputs under the stated semantics, whereas tests can expose
implementation discrepancies only on the inputs that they exercise.

\subsection{Structure, traces and branching are non-interchangeable}
\label{sec:discussion-layers}

Proposition \ref{prop:label-blind} gives a general reason why identical
unlabelled structure cannot determine labelled behaviour. The controlled label
mismatch and order swap instantiate that limitation: both receive a PN-GDDA
score of 1.0 although neither simulation direction holds. Example
\ref{ex:branching} supplies the complementary linear-time limitation. Its two
systems have the same exact finite traces, yet they differ in the point at which
choice is resolved and satisfy only one simulation direction.

Direct PN-GDDA achieved 2/6 on the predeclared binary task, but this is not a
software failure or a claim that PN-GDDA is unsuitable for its intended purpose.
The independent Holmes score agrees to 12 decimal places, and the 592- and
576-slot catalogues give the same threshold decisions. Transition labels,
initial marking and reachable branching lie outside the structural statistic by
design. LTS-GDA recovers labelled local motifs and improves to 4/6, but it still
does not implement silent abstraction or recursive successor matching. Its
held-out association is the strongest observed ($\rho=0.669$), while the exact
$p=0.097$ remains inconclusive. The appropriate conclusion is to report
structural, trace and branching views as answers to different questions.

\subsection{Observational interfaces and abstraction}
\label{sec:discussion-interface}

The relation is indexed by $\Sigma$. Relabelling an event as internal changes
which paths are abstracted, and merging distinct outcomes under one label can
coarsen the result. In the GIM application, apoptosis and
differentiation or endoreduplication share the declared observable
\texttt{exit\_cycle}. Weak bisimilarity therefore describes agreement only at
that selected resolution. It is not a statement that the terminal mechanisms are
biologically identical.

Proposition \ref{prop:silent-refinement} justifies the controlled subdivisions
used in the robustness analysis. It does not license arbitrary Petri-net
refinement: added conflicts, interleavings or observable choices may change the
reachable LTS. Update semantics is another part of the model contract. The
asynchronous--synchronous GINsim differences and the two changed HPN classes show
that a relation should not be described as semantics-neutral. Most-permissive,
general-asynchronous and stochastic conventions were not evaluated
\cite{Pauleve2020}.

\subsection{Independent computational evidence}
\label{sec:discussion-validation}

The six synthetic pairs isolate known relations before any biological case is
examined. The formal classifier recovers all six, while the comparator errors
occur in the cases predicted by their information loss. Across the six synthetic
pairs, six public-model controls and nine HPN comparisons, Python and mCRL2 agree
on all 42 reported strong/weak decisions. mCRL2 additionally confirms the exact
weak-trace controls, including acceptance of the branching pair that both
bisimulation implementations reject.

The dual attacker--defender implementation checks the simulation orientation not
available through the used \texttt{ltscompare} interface. Its agreement on all
67,600 ordered pairs is exhaustive only within the declared universe of 260
one- and two-state LTSs over $\{a,\TAU\}$. It is evidence against small-system
implementation defects, not a proof for arbitrary systems. The transformed
GINsim controls similarly test parsing, execution and export, but their expected
relations follow from controlled transformations. Distinct HPN model families
and held-out measurements provide a separate empirical check rather than a
replacement for formal validation.

\subsection{Biological applications and empirical limits}
\label{sec:discussion-biological}

The HPN-DREAM application has nine pair--condition observations from three cell
lines. Six comparisons show one-way simulation and three show
non-comparability; none is weakly bisimilar. The ordinal formal class has
$\rho=0.091$ with experimental distance and an inconclusive exact $p=0.111$.
Truncated trace distance also lacks supportive concordance. CASPOTS reproduces
held-out compatibility, but the native-versus-cross advantage is not established
($p=0.25$). CASPOTS RMSE is a minimum-distance compatibility score over
admissible trajectories, not a unique time-series prediction \cite{Razzaq2018}.
These results do not support prognostic use or cell-line specificity.

The curated Petri-net outputs are also conditional. The models omit kinetics,
quantities, noise, stoichiometric uncertainty, tissue context and independent
experimental inference. Different curators could select different arcs, initial
markings or observables. Provenance coverage documents the choices but does not
make them independent evidence. Label-permutation values measure sensitivity
inside the constructed models and are not biological significance tests. In
particular, weak relations between an Arabidopsis net and an animal net do not
establish organism-level functional equivalence.

\subsection{Computational limitations and use}
\label{sec:discussion-computation}

The implementation enumerates states explicitly and stores up to
$|S_1||S_2|$ candidate pairs. Silent closures are cached, but visible weak target
sets are reconstructed on demand. The scaling study reaches 145 states and
18,705 candidate pairs; it does not establish genome-scale performance. Symbolic
state representations, partition refinement and compositional checking are
appropriate directions for larger systems.

PN-GDDA is evaluated directly on the native curated nets and canonical
state-machine encodings of the six controlled LTS pairs. Applying it to
SBML-qual or inferred Boolean rules would require a non-unique reverse translation
into Petri-net structure. The difference between 592 published/Holmes slots and
576 automorphism-distinct orbits also remains a catalogue question beyond the
compact sensitivity analysis. Future evaluations should use larger repository
Petri nets, prospectively expert-labelled relations, biological replicates,
additional update semantics and blinded interfaces. In its current form, the
framework is most appropriate for auditing two explicit qualitative models at a
predeclared observational resolution.

\section{Conclusion}\label{sec:conclusion}

We formulated a graded comparison framework for finite labelled Petri nets and
other executable qualitative models under a declared interface. The classifier
reports the strongest supported relation among strong bisimilarity, weak
bisimilarity, mutual or one-way weak simulation and non-comparability. Its
priority follows from the relation hierarchy, and the implemented deletion
algorithm computes the corresponding greatest fixed points and terminates on
finite LTSs. Controlled silent refinement, label relabelling and the
trace-equivalent branching example establish why internal abstraction,
unlabelled structure, linear traces and branching cannot be treated as
interchangeable.

Construction-level cases, mCRL2 and the independent simulation game provide
separate computational checks, including complete agreement on 67,600 ordered
pairs within the declared small-LTS universe. Applications to GINsim,
HPN-DREAM/CASPOTS and curated Petri nets demonstrate cross-format execution and
make interface and update-semantics dependence visible. They do not establish
predictive performance, cell-line specificity or organism-level equivalence.
The resulting contribution is a reproducible formal model audit whose
mathematical outputs and empirical diagnostics remain explicitly separated.

\section*{Data and code availability}

Curated definitions, hash-verified public GINsim models, HPN-DREAM/CASPOTS model
families, normalized perturbation-response files, benchmark metadata and generated
result tables are available at
\url{https://github.com/cero1979/BisimulationMol}. Source code is released under
the MIT License. The repository documents separate commands for external-model
validation, HPN-DREAM validation, the Holmes PN-GDDA check, figure regeneration,
unit tests and notebook execution. No confidential or participant-level data were
used.

\section*{Funding}

The author received no specific funding for this work.

\section*{Competing interests}

The author declares no financial or non-financial interests directly or
indirectly related to this work.

\section*{Ethics statement}

Not applicable. The study uses computational models and previously published
aggregate information; no human participants, human samples or live animals were
involved.

\section*{Author contributions}

Carlos Ramirez Ovalle: conceptualization, methodology, software, formal analysis,
validation and visualization; writing---original draft; writing---review and
editing.

\section*{Declaration of AI Use}

During the preparation of this work, the author used OpenAI Codex for the
following purposes: improving language and readability, checking grammar,
formatting references, adapting the LaTeX source to the journal template, and
assisting with reproducibility checks. After using this tool, the author reviewed
and edited the content as needed and takes full responsibility for the content of
the publication.
% Alphabetical references in the format shown on the JMCS Journal Format page.
% Metadata and DOI resolution are documented in REFERENCE_AUDIT.md.
\begin{thebibliography}{99}

\bibitem{AbouJaoude2009}
\href{https://doi.org/10.1016/j.jtbi.2009.02.005}
{W. Abou-Jaoud\'e, D. A. Ouattara, M. Kaufman,
\textit{From structure to dynamics: Frequency tuning in the p53--Mdm2 network},
J. Theor. Biol., {\bf 258} (2009), 561--577.}

\bibitem{BaierKatoen2008}
C. Baier, J.-P. Katoen, \textit{Principles of Model Checking}, MIT Press,
Cambridge, MA, (2008).

\bibitem{Benes2023}
\href{https://doi.org/10.1093/bioinformatics/btad158}
{N. Bene{\v{s}}, L. Brim, O. Huvar, S. Pastva, D. {\v{S}}afr\'anek,
\textit{Boolean network sketches: A unifying framework for logical model
inference}, Bioinformatics, {\bf 39} (2023), btad158.}

\bibitem{Bunte2019}
\href{https://doi.org/10.1007/978-3-030-17465-1_2}
{O. Bunte, J. F. Groote, J. J. A. Keiren, M. Laveaux, T. Neele, E. P. de Vink,
W. Wesselink, A. Wijs, T. A. C. Willemse, \textit{The mCRL2 toolset for
analysing concurrent systems: Improvements in expressivity and usability}, in:
Tools and Algorithms for the Construction and Analysis of Systems, LNCS
{\bf 11428}, Springer, (2019), 21--39.}

\bibitem{ChaouiyaPetriBio2007}
\href{https://doi.org/10.1093/bib/bbm029}
{C. Chaouiya, \textit{Petri net modelling of biological networks}, Brief.
Bioinform., {\bf 8} (2007), 210--219.}

\bibitem{Chaouiya2013SBML}
\href{https://doi.org/10.1186/1752-0509-7-135}
{C. Chaouiya, D. B\'erenguier, S. M. Keating, A. Naldi, M. P. van Iersel,
N. Rodriguez, et al., \textit{SBML qualitative models: A model representation
format and infrastructure to foster interactions between qualitative modelling
formalisms and tools}, BMC Syst. Biol., {\bf 7} (2013), 135.}

\bibitem{CifuentesFontanals2022}
\href{https://doi.org/10.3389/fams.2022.838546}
{L. Cifuentes-Fontanals, E. Tonello, H. Siebert, \textit{Control in Boolean
networks with model checking}, Front. Appl. Math. Stat., {\bf 8} (2022), 838546.}

\bibitem{ClavijoBuritica2023}
\href{https://doi.org/10.1016/j.heliyon.2023.e15367}
{D. C. Clavijo-Buritic\'a, C. C. Sosa, R. C\'ardenas Heredia, A. J. Mosquera,
A. \'Alvarez, J. Medina, M. Quimbaya, \textit{Use of Arabidopsis thaliana as a
model to understand specific carcinogenic events: Comparison of the molecular
machinery associated with cancer-hallmarks in plants and humans}, Heliyon,
{\bf 9} (2023), e15367.}

\bibitem{DeVeylder2007}
\href{https://doi.org/10.1038/nrm2227}
{L. De Veylder, T. Beeckman, D. Inz\'e, \textit{The ins and outs of the plant
cell cycle}, Nat. Rev. Mol. Cell Biol., {\bf 8} (2007), 655--665.}

\bibitem{Fisher2007}
\href{https://doi.org/10.1038/nbt1356}
{J. Fisher, T. A. Henzinger, \textit{Executable cell biology}, Nat.
Biotechnol., {\bf 25} (2007), 1239--1249.}

\bibitem{Hagberg2008}
A. A. Hagberg, D. A. Schult, P. J. Swart, \textit{Exploring network structure,
dynamics, and function using NetworkX}, in: Proceedings of the 7th Python in
Science Conference, Pasadena, CA, (2008), 11--15.

\bibitem{Hanahan2011}
\href{https://doi.org/10.1016/j.cell.2011.02.013}
{D. Hanahan, R. A. Weinberg, \textit{Hallmarks of cancer: The next generation},
Cell, {\bf 144} (2011), 646--674.}

\bibitem{Hanahan2022}
\href{https://doi.org/10.1158/2159-8290.CD-21-1059}
{D. Hanahan, \textit{Hallmarks of cancer: New dimensions}, Cancer Discov.,
{\bf 12} (2022), 31--46.}

\bibitem{Heiner2008}
\href{https://doi.org/10.1007/978-3-540-68894-5_7}
{M. Heiner, D. Gilbert, R. Donaldson, \textit{Petri nets for systems and
synthetic biology}, in: Formal Methods for Computational Systems Biology, LNCS
{\bf 5016}, Springer, Berlin, Heidelberg, (2008), 215--264.}

\bibitem{Hill2016}
\href{https://doi.org/10.1038/nmeth.3773}
{S. M. Hill, L. M. Heiser, T. Cokelaer, M. Unger, N. K. Nesser, D. E. Carlin,
et al., \textit{Inferring causal molecular networks: Empirical assessment
through a community-based effort}, Nat. Methods, {\bf 13} (2016), 310--318.}

\bibitem{Koch2005}
\href{https://doi.org/10.1093/bioinformatics/bti145}
{I. Koch, B. H. Junker, M. Heiner, \textit{Application of Petri net theory for
modelling and validation of the sucrose breakdown pathway in the potato tuber},
Bioinformatics, {\bf 21} (2005), 1219--1226.}

\bibitem{Marshall2018}
\href{https://doi.org/10.1146/annurev-arplant-042817-040606}
{R. S. Marshall, R. D. Vierstra, \textit{Autophagy: The master of bulk and
selective recycling}, Annu. Rev. Plant Biol., {\bf 69} (2018), 173--208.}

\bibitem{Milner1989}
R. Milner, \textit{Communication and Concurrency}, Prentice Hall, New York,
(1989).

\bibitem{Murata1989}
\href{https://doi.org/10.1109/5.24143}
{T. Murata, \textit{Petri nets: Properties, analysis and applications}, Proc.
IEEE, {\bf 77} (1989), 541--580.}

\bibitem{Naldi2018}
\href{https://doi.org/10.3389/fphys.2018.00646}
{A. Naldi, C. Hernandez, W. Abou-Jaoud\'e, P. T. Monteiro, C. Chaouiya,
D. Thieffry, \textit{Logical modeling and analysis of cellular regulatory
networks with GINsim 3.0}, Front. Physiol., {\bf 9} (2018), 646.}

\bibitem{Pauleve2020}
\href{https://doi.org/10.1038/s41467-020-18112-5}
{L. Paulev\'e, J. Kol{\v{c}}\'ak, T. Chatain, S. Haar, \textit{Reconciling
qualitative, abstract, and scalable modeling of biological networks}, Nat.
Commun., {\bf 11} (2020), 4256.}

\bibitem{Quimbaya2012}
\href{https://doi.org/10.1007/s00018-011-0909-x}
{M. Quimbaya, K. Vandepoele, E. Rasp\'e, M. Matthijs, S. Dhondt,
G. T. S. Beemster, G. Berx, L. De Veylder, \textit{Identification of putative
cancer genes through data integration and comparative genomics between plants
and humans}, Cell. Mol. Life Sci., {\bf 69} (2012), 2041--2055.}

\bibitem{Radom2017}
\href{https://doi.org/10.1093/bioinformatics/btx492}
{M. Radom, A. Rybarczyk, B. Szawulak, H. Andrzejewski, P. Chabelski, A. Kozak,
P. Formanowicz, \textit{Holmes: A graphical tool for development, simulation
and analysis of Petri net based models of complex biological systems},
Bioinformatics, {\bf 33} (2017), 3822--3823.}

\bibitem{Razzaq2018}
\href{https://doi.org/10.1371/journal.pcbi.1006538}
{M. Razzaq, L. Paulev\'e, A. Siegel, J. Saez-Rodriguez, J. Bourdon,
C. Guziolowski, \textit{Computational discovery of dynamic cell line specific
Boolean networks from multiplex time-course data}, PLOS Comput. Biol., {\bf 14}
(2018), e1006538.}

\bibitem{Regev2002}
\href{https://doi.org/10.1038/419343a}
{A. Regev, E. Shapiro, \textit{Cellular abstractions: Cells as computation},
Nature, {\bf 419} (2002), 343.}

\bibitem{Sangiorgi2011}
\href{https://doi.org/10.1017/CBO9780511777110}
{D. Sangiorgi, \textit{Introduction to Bisimulation and Coinduction}, Cambridge
University Press, Cambridge, (2011).}

\bibitem{Szawulak2022}
\href{https://doi.org/10.1038/s41598-022-24535-5}
{B. Szawulak, P. Formanowicz, \textit{Graphlets in comparison of Petri
net-based models of biological systems}, Sci. Rep., {\bf 12} (2022), 20942.}

\bibitem{Traynard2016}
\href{https://doi.org/10.1093/bioinformatics/btw457}
{P. Traynard, A. Faur\'e, F. Fages, D. Thieffry, \textit{Logical model
specification aided by model-checking techniques: Application to the mammalian
cell cycle regulation}, Bioinformatics, {\bf 32} (2016), i772--i780.}

\bibitem{Yoshiyama2013}
\href{https://doi.org/10.3390/biology2041338}
{K. O. Yoshiyama, K. Sakaguchi, S. Kimura, \textit{DNA damage response in
plants: Conserved and variable response compared to animals}, Biology, {\bf 2}
(2013), 1338--1356.}

\end{thebibliography}


\end{document}
